% !Mode:: "TeX:UTF-8"
% !TEX encoding = UTF-8 Unicode 
% vim: set fenc=utf-8 ft=latex
% -*- mode: latex -*-
% -*- coding: UTF-8 -*-
% kate: encoding utf-8;
\documentclass{stagetex}
\usepackage[utf8]{inputenc}


% Qui dovresti vedere lettere accentate: èéìàòù

%Inserisci qui il tuo nome
\Stagista{Pinco Pallino}
\Stage{Senior 2014}
\Scuola{Liceo Speciale di Interesse Olimpico ``D. Schwarz'' -- Cesenatico}

%Inserisci qui il tipo di esercizi che devi svolgere -- vedi istruzioni sul forum per sapere qual è
\Esercizi{PreIMO -- P}

%Per gli stage Senior, si intende l'anno che state per *iniziare*, non quello che avete già finito
%Il primo anno delle superiori è 1, il quinto (quello della maturità per la maggior parte degli studenti) è 5. Attenetevi a questa numerazione
%anche se frequentate un liceo classico.
\AnnoCorso{5}

\Genere{M} % maschio, femmina, oppure anche ``è complicato''. Scusate le informazioni sensibili, ma ci è utile per dividervi nelle camere di albergo.

\begin{document}

\begin{soluzione}{A1}
 Scrivete le vostre soluzioni tra questi due comandi. Una nuova pagina viene creata automaticamente per ogni esercizio.
\end{soluzione}

\begin{soluzione}{A2}
Per creare un nuovo paragrafo, lasciate una linea vuota (due a capo successivi).
 
Non serve ricopiare il testo dell'esercizio. 
Scrivete in maniera \emph{sintetica}, ma \emph{chiara} e \emph{completa}
tutti i passaggi necessari. Non bisogna limitarsi a ricopiare l'aiutino o a
sbobinare il video, fermandosi con dei puntini dove il video dice ``e da
qui si conclude con facili calcoli''. Se ci sono delle verifiche da fare, fatele!

Si possono inserire formule tra i dollari, per esempio, sia $a\in\mathbb{R}$, $a^5>0$. Oppure formule che prendono una riga intera,
come
\begin{equation}
    e^{i\pi}+1=0
\end{equation}
Esponenti o pedici più grossi di una lettera vanno messi tra parentesi graffe: $a^5 > a^{10}$.

Potete dare un'etichetta (\emph{label}) alle formule e riferirvi ad esse più avanti nel testo quando vi serve. Per esempio, consideriamo il \textbf{teorema di Pitagora}
\begin{equation}
 a^2+b^2=c^2, \label{pitagora}
\end{equation}
insieme con la nota formula
\begin{equation}
E=mc^2. \label{Einstein}
\end{equation}
Combinando la \eqref{pitagora} e la \eqref{Einstein}, segue banalmente
\begin{equation}
 E=m(a^2+b^2).
\end{equation}

\end{soluzione}

\begin{soluzione}{G4}
Dotatevi di un programma per trasformare file \texttt{.tex} in file \texttt{.pdf}, per esempio
 \begin{itemize}
  \item per Windows, MikTex, che potete recuperare su \url{http://miktex.org/};
  \item per Mac OS, BasicTeX \url{https://www.tug.org/mactex/morepackages.html} oppure MacTex \url{https://www.tug.org/mactex/} (grosso da scaricare)
  \item per Linux, Texlive (su Debian o Ubuntu, installate il pacchetto \texttt{texlive} tramite usando il Software Center).
 \end{itemize}
Se non volete scaricare nulla, potete anche provare ad usare ShareLatex \url{https://www.sharelatex.com/}, che vi permette di compilare file .tex online.
In ogni caso, impostate il vostro programma in modo da produrre direttamente un file pdf (tex=>pdf in MikTeX).

Alcuni di questi programmi sono già distribuiti con un editor abbastanza sofisticato per scrivere file \texttt{.tex}, come Texmaker \url{http://www.xm1math.net/texmaker/}.
In ogni caso si tratta di normali file di testo, quindi va benissimo anche il Blocco Note o un editor generico come Notepad++ \url{http://notepad-plus-plus.org/}.

Ci sono molte diverse configurazioni possibili e non ho provato tutto; fatemi avere del feedback attraverso il forum su cosa funziona meglio/peggio.

Perché tutto funzioni, i due file \texttt{template.tex} (o come l'avete rinominato) e \texttt{stagetex.cls} (che deve restare con questo nome)
devono stare nella stessa cartella, così come eventuali immagini.

Alcune volte è necessario ricompilare una seconda volta perché tutti i cambiamenti nella numerazione siano riconosciuti.

I messaggi di errore di \LaTeX{} sono universalmente considerati orribili. Potete perdere anche mezz'ora o un'ora a causa un errore incomprensibile. :(
La sezione ``LaTeX, questo sconosciuto'' del forum e \url{http://tex.stackexchange.com/} sono vostri amici. Buona fortuna!
\end{soluzione}

\begin{soluzione}{G7}
 Potete inserire immagini con i comandi che vedete qui sotto.
 \begin{center}
 \includegraphics[width=0.9\textwidth]{fanohere.pdf} 
 \end{center}
Dovrebbero funzionare file con estensione \texttt{.pdf}, \texttt{.png} o \texttt{.jpg},
da piazzare nella stessa cartella del documento.
Il primo è preferibile perché genera immagini di qualità maggiore (non sgranate e pixellose).

L'opzione \texttt{width} serve per indicare la larghezza (in multipli della larghezza del testo \texttt{textwidth}.

Un ottimo tool per creare figure di geometria da inserire nelle soluzioni è Geogebra, \url{http://www.geogebra.org/}. Altrimenti vanno benissimo anche disegni fatti a mano e scannerizzati.
\end{soluzione}

\begin{soluzione}{A8}
Ecco qualche trucco in più per creare cose più complicate nelle vostre soluzioni!

\begin{lemma} \label{ilmiolemma}
 Per ogni numero reale $a$, si ha che $a=2a$.
\end{lemma}
\begin{proof}
 Chiaramente $1=2$. Quindi, moltiplicando per $a$, $a=2a$.
\end{proof}

Utilizzando il Lemma \ref{ilmiolemma}, concludiamo!

Equazione più complicata
\begin{equation}
 \sum_{k=0}^{\infty} a^k = \frac{1}{1-a}  \qquad \text{per ogni $a$ con $\abs{a}<1$.}
\end{equation}
Raccogliendo, abbiamo
\begin{equation}
 \underbrace{1+a+a^2+\dots+a^k}_{ = \frac{1-a^{k+1}}{1-a}} \geq b^3-3b+37.
\end{equation}
Equazioni su più righe:
\begin{align}
 1-\frac{1}{n+1} &= \left(1-\frac{1}{2}\right) + \left(\frac{1}{2} - \frac{1}{3}\right)+ \dots + \left(\frac{1}{n} - \frac{1}{n+1}\right)\\
                 &= \frac{2-1}{1\cdot 2} + \frac{3-2}{2\cdot 3} + \dots + \frac{n+1-n}{n(n+1)}\\
                 &= \frac{1}{1\cdot 2} + \frac{1}{2\cdot 3} + \dots + \frac{1}{n(n+1)}\\
                 &= \sum_{k=1}^n \frac{1}{k(k+1)}.
\end{align}
Disuguaglianze! Si ha
\begin{equation}
 ab+bc+ca \stackrel{(*)}{\leq} a^2+b^2+c^2 \stackrel{?}{\leq} 1,
\end{equation}
dove la $(*)$ è vera per riarrangiamento. O bunching. O Cauchy-Schwarz. Insomma, non può non venirvi.
\end{soluzione}

\begin{soluzione}{M5}
 \begin{steplist}
  \step{Uguaglianza tra gli angoli} Prima di tutto, dimostriamo che $\widehat{ABC}=\angle CAB$. Quindi $\triangle ABC$ è isoscele;
  sia il suo cerchio inscritto $\odot ABC$.
  \step{Uguaglianza tra i cerchi} Preriscaldate il forno a $220^\circ$, poi infornate insieme al gatto di casa. La torta cuoce in 50 minuti, il gatto no.
  \step{Concludiamo} Il resto è trivial.
 \end{steplist}
 
 Se avete bisogno di trovare i comandi per inserire simboli particolari, potete usare \url{http://detexify.kirelabs.org/},
 che riconosce (o almeno ci prova\dots) un simbolo disegnato a mano, oppure cercarlo sulla (lunga) lista su
 \url{http://mirrors.ctan.org/info/symbols/comprehensive/symbols-a4.pdf}.
\end{soluzione}
\begin{soluzione}{M6}
  Buona fortuna nella scrittura delle soluzioni, e speriamo di vederci a Pisa!
\end{soluzione}

\end{document}
